\documentclass[11pt,a4paper]{article}
\usepackage[margin=2.2cm]{geometry}
\usepackage{enumitem}
\usepackage{booktabs}
\usepackage{array}
\usepackage[hidelinks]{hyperref}
\setlist{itemsep=1pt,topsep=3pt}
\setlength{\parindent}{0pt}
\setlength{\parskip}{4pt}

\begin{document}

\begin{center}
{\LARGE\bfseries Lightweight Temperature Predictor:\\[2pt] Calibrate Once, Predict Fast}\\[6pt]
{\large Scope Document --- Computer Architecture Course Project}\\[2pt]
Topic: \emph{Thermal management: lightweight temperature predictor}
\end{center}

\section*{1. Team}
\begin{tabular}{@{}ll@{}}
Ethan Bobby Kurien & HotSpot tooling, step-response calibration, closed-loop DTM setup\\
Kaivalya Mase & floorplans, trace-based calibration, forecasting, DTM policies\\
\end{tabular}

\section*{2. Problem statement}
For a given heterogeneous multicore floorplan, build a predictor that returns steady-state or
transient block temperatures for any power input, much faster than HotSpot and with a very small
model. The predictor works in two stages: a one-time \emph{calibration} per floorplan, which runs
HotSpot a small number of times to fit a compact set of coefficients, and a \emph{prediction}
stage that uses only those coefficients. HotSpot serves as the reference. Performance is judged on
\textbf{accuracy}, \textbf{prediction time} and \textbf{model size}.

\section*{3. End deliverables}
\subsection*{Must-have}
\begin{itemize}
  \item A HotSpot-based ground-truth setup: run and parse scripts, and a few heterogeneous
        floorplans with matching power ranges.
  \item A calibration stage that fits predictor coefficients for a given floorplan.
  \item A steady-state predictor.
  \item A transient predictor.
  \item Evaluation of both predictors against HotSpot on accuracy, prediction time and model size,
        including a check that per-floorplan calibration is actually needed.
  \item A small Python tool packaging calibration and prediction.
\end{itemize}
\subsection*{Stretch goals}
\begin{itemize}
  \item Proactive DTM: using the predictor's forecast to throttle or migrate tasks before a
        temperature limit is crossed, compared against reactive DTM.
  \item Robustness studies (unseen workloads, package changes, leakage).
  \item Model-size reduction for larger core counts.
\end{itemize}

\section*{4. Plan and approach}
\begin{enumerate}
  \item \textbf{Ground truth:} establish HotSpot as a trusted, scriptable reference and confirm
        that it behaves linearly without leakage, which is what makes a compact predictor possible.
  \item \textbf{Calibration:} fit coefficients from a small number of planned HotSpot runs,
        keeping calibration inputs separate from test inputs.
  \item \textbf{Prediction and evaluation:} implement steady and transient predictors,
        compare against simple baselines, and measure the three metrics.
  \item \textbf{Extensions:} package the tool, then attempt the DTM stretch goal as time permits.
\end{enumerate}
Work is split between the two members along the calibration approach each owns, with shared file
formats agreed early and all work tracked on GitHub.

\section*{5. Anticipated risks}
\begin{itemize}
  \item \textbf{No given floorplan or real workload power:} synthetic heterogeneous floorplans and
        power traces will be used; the method is not tied to a particular floorplan.
  \item \textbf{HotSpot output precision (0.01\,K rounding):} to be addressed by patching the
        output format.
  \item \textbf{Block-level model may hide hot spots inside large cores:} to be quantified
        against HotSpot's grid model.
  \item \textbf{Stretch goals may not fit in the available time:} they are ordered last and
        can be dropped without affecting the core deliverables.
\end{itemize}

\section*{6. LLM assistance}
Claude (Anthropic) was used to help draft the project plan, literature review, a floorplan
generation script, and this document. All output was reviewed by the team. Any further use will be
disclosed in the same way.

\end{document}
